Grundsätzlich wird jede Seite einer Arbeit nummeriert. Die Nummierung beginnt
mit dem Deckblatt, bei Buchveröffentlichungen mit der ersten Seite der Titelei.
Der Buchumschlag wird nicht nummieriert. Auf bestimmten Seiten (z.\,B. dem
Deckblatt oder auf leeren Seiten) wird keine Seitennummer gedruckt. Bei
umfangreicheren Arbeiten, beginnt die Nummierierung ab dem ersten Kapitel bzw.
dem Vorwort erneut bei~1. Die davor liegenden Seiten (Inhaltsverzeichnis etc.),
werden dann der Unterscheidbarkeit halber mit römischen Zahlen nummieriert.
Die Umschaltung der Nummierierung erfolgt mit dem \LaTeX-Befehl
\verb+\pagenumbering+, wie in Listing~\ref{lst:final:pagenumber} dargestellt.
Der Seitenzähler wird automatisch auf~1 zurückgesetzt.

\lstinputlisting[language={[LaTeX]{TeX}},style=block,escapechar=\!,float={htb},caption={Umschalten der Seitennummerierung},label={lst:final:pagenumber}]{code/final_pagenumber.tex}

In der Dokumentklasse~\texttt{book} erfolgt der Wechsel der Nummerierung
automatisch durch die Kommandos \verb+\frontmatter+, \verb+\mainmatter+ und
\verb+\backmatter+.

Ein separates Nummerierungsschema für Anhänge ist normalerweise unüblich.
Sollen Anhänge in der Form A-1, A-2, \dots, B-1, B-2, \dots\ durchnummierert
werden, muss gemäß Listing~\ref{lst:final:pagenumber_appendix} verfahren werden.
Das Konstrukt aus \verb+\pagenumbering+ und \verb+\renewcommand+ muss mit jedem
Anhang wiederholt werden. Entsprechende Anpassungen sind dann auch für die
folgenden Verzeichnisse, wie das Literaturverzeichnung, Glossar, etc.
vorzunehmen. Eine durchlaufende arabische Nummerierung ist hier die elegantere
Lösung.

\lstinputlisting[language={[LaTeX]{TeX}},style=block,escapechar=\!,float={htb},caption={Separate Seitennummerierung für Anhänge},label={lst:final:pagenumber_appendix},morekeywords={chapter}]{code/final_pagenumber_appendix.tex}
